% GENERATED FILE. Edit canonical lessons, then run npm run generate:books.
% canonical-source-hash: fnv1a64:b202cc876c25077a
% canonical-lessons: MR-C458-dilasa, MR-C458-dhir, MR-C458-santvan-karne, MR-C458-hayse-vatne, MR-C458-gahivarne

\chapter{Comfort and Relief}
\label{ch:mr-comfort-and-relief}
\hlchaptermodality{458}

\begin{chapteropening}
\textbf{By the end of this chapter:} \emph{I can comfort someone and say that I feel relieved.}

Complete the last lesson of chapter 458: I can comfort someone and say that I feel relieved.
\end{chapteropening}

\section[\texorpdfstring{\mr{दिलासा} (dilāsā)}{दिलासा (dilāsā)}]{\mr{दिलासा} (dilāsā) — comfort; reassurance}

\label{lesson:MR-C458-dilasa}



\begin{quote}
Before the new one: say the Marathi for to feel sorry, then the Marathi for to sob quietly.
\end{quote}

\subsection*{You'll want to know: \mr{दिलासा}}
\textbf{\mr{दिलासा}} — \emph{dilāsā} — \textquotedblleft{}comfort; reassurance\textquotedblright{}.

\begin{quote}
\mr{मित्राच्या} \mr{फोनने} \mr{मला} \mr{खूप} \mr{दिलासा} \mr{मिळाला}.
\end{quote}

\begin{quote}
\emph{mitrācyā phonne malā khūp dilāsā miḷālā.}
\end{quote}

\begin{quote}
My friend's phone call was a great comfort to me.
\end{quote}

Masculine.

\mr{दिलासा} \mr{देणे}, to reassure; \mr{दिलासा} \mr{मिळणे}, to be reassured, with the dative.

\begin{cousinweb}
\textbf{\mr{दिलासा}} (\emph{dilāsā}) is a loan from Persian \emph{dilāsā}, \textquotedblleft{}consolation\textquotedblright{}, built on \emph{dil}, \textquotedblleft{}heart\textquotedblright{}.
\end{cousinweb}

\begin{grammarlens}[title={What you've built}]
One more word for describing what you lived and felt.
\end{grammarlens}

\subsection*{Your turn}
\begin{itemize}
\raggedright
  \item \emph{Say it:} \emph{dilāsā}
  \item \emph{Say it:} \emph{dilāsā}, once more
  \item \emph{Say it:} \emph{mitrācyā phonne malā khūp dilāsā miḷālā.} — My friend's phone call was a great comfort to me.
  \item \emph{Recall:} say the Marathi for to feel sorry, then the Marathi for to sob quietly, then say \emph{dilāsā} again
\end{itemize}

\begin{tcolorbox}[breakable,colback=teal!4,colframe=teal!35!black,title={Before you move on}]
What does \mr{दिलासा} mean? (\textquotedblleft{}Comfort; reassurance\textquotedblright{}.) Say it once more.
\end{tcolorbox}
\section[\texorpdfstring{\mr{धीर} (dhīr)}{धीर (dhīr)}]{\mr{धीर} (dhīr) — fortitude; heart}

\label{lesson:MR-C458-dhir}



\begin{quote}
Before the new one: say the Marathi for to sob quietly, then the Marathi for comfort.
\end{quote}

\subsection*{You'll want to know: \mr{धीर}}
\textbf{\mr{धीर}} — \emph{dhīr} — \textquotedblleft{}fortitude; heart\textquotedblright{}.

\begin{quote}
\mr{आईने} \mr{धीर} \mr{दिला} \mr{म्हणून} \mr{मी} \mr{पुन्हा} \mr{प्रयत्न} \mr{केला}.
\end{quote}

\begin{quote}
\emph{āīne dhīr dilā mhaṇūn mī punhā prayatna kelā.}
\end{quote}

\begin{quote}
Mother gave me heart, so I tried again.
\end{quote}

Masculine.

\mr{धीर} \mr{देणे}, to give someone heart; \mr{धीर} \mr{धरणे}, to be patient.

\begin{cousinweb}
\textbf{\mr{धीर}} (\emph{dhīr}) is a tatsama, Sanskrit \emph{dhīra}, \textquotedblleft{}wise; steady, resolute\textquotedblright{}; CDIAL lists Marathi \emph{dhīr}, \textquotedblleft{}patience\textquotedblright{} (CDIAL 6817).
\end{cousinweb}

\begin{grammarlens}[title={What you've built}]
One more word for describing what you lived and felt.
\end{grammarlens}

\subsection*{Your turn}
\begin{itemize}
\raggedright
  \item \emph{Say it:} \emph{dhīr}
  \item \emph{Say it:} \emph{dhīr}, once more
  \item \emph{Say it:} \emph{āīne dhīr dilā mhaṇūn mī punhā prayatna kelā.} — Mother gave me heart, so I tried again.
  \item \emph{Recall:} say the Marathi for to sob quietly, then the Marathi for comfort, then say \emph{dhīr} again
\end{itemize}

\begin{tcolorbox}[breakable,colback=teal!4,colframe=teal!35!black,title={Before you move on}]
What does \mr{धीर} mean? (\textquotedblleft{}Fortitude; heart\textquotedblright{}.) Say it once more.
\end{tcolorbox}
\section[\texorpdfstring{\mr{सांत्वन} \mr{करणे} (sāntvan karṇe)}{सांत्वन करणे (sāntvan karṇe)}]{\mr{सांत्वन} \mr{करणे} (sāntvan karṇe) — to console}

\label{lesson:MR-C458-santvan-karne}



\begin{quote}
Before the new one: say the Marathi for comfort, then the Marathi for fortitude.
\end{quote}

\subsection*{You'll want to know: \mr{सांत्वन} \mr{करणे}}
\textbf{\mr{सांत्वन} \mr{करणे}} — \emph{sāntvan karṇe} — \textquotedblleft{}to console\textquotedblright{}.

\begin{quote}
\mr{मित्रांनी} \mr{आजीचं} \mr{सांत्वन} \mr{केलं}.
\end{quote}

\begin{quote}
\emph{mitrānnī ājīcaṁ sāntvan kelaṁ.}
\end{quote}

\begin{quote}
Friends consoled Grandmother.
\end{quote}

Past, with -\mr{ने} (he …): \mr{त्याने} \mr{तिचं} \mr{सांत्वन} \mr{केलं}, \emph{tyāne ticaṁ sāntvan kelaṁ}.

The one consoled takes the genitive: \mr{मित्राचं} \mr{सांत्वन} \mr{करणे}.

\begin{cousinweb}
\textbf{\mr{सांत्वन} \mr{करणे}} (\emph{sāntvan karṇe}) is the tatsama \mr{सांत्वन}, Sanskrit \emph{sāntvana}, \textquotedblleft{}soothing with kind words\textquotedblright{} (Monier-Williams), with the light verb \mr{करणे}.
\end{cousinweb}

\begin{grammarlens}[title={What you've built}]
One more word for describing what you lived and felt.
\end{grammarlens}

\subsection*{Your turn}
\begin{itemize}
\raggedright
  \item \emph{Say it:} \emph{sāntvan karṇe}
  \item \emph{Say it:} \emph{sāntvan karṇe}, once more
  \item \emph{Say it:} \emph{mitrānnī ājīcaṁ sāntvan kelaṁ.} — Friends consoled Grandmother.
  \item \emph{Recall:} say the Marathi for comfort, then the Marathi for fortitude, then say \emph{sāntvan karṇe} again
\end{itemize}

\begin{tcolorbox}[breakable,colback=teal!4,colframe=teal!35!black,title={Before you move on}]
What does \mr{सांत्वन} \mr{करणे} mean? (\textquotedblleft{}To console\textquotedblright{}.) Say it once more.
\end{tcolorbox}
\section[\texorpdfstring{\mr{हायसे} \mr{वाटणे} (hāyse vāṭṇe)}{हायसे वाटणे (hāyse vāṭṇe)}]{\mr{हायसे} \mr{वाटणे} (hāyse vāṭṇe) — to feel relieved}

\label{lesson:MR-C458-hayse-vatne}



\begin{quote}
Before the new one: say the Marathi for fortitude, then the Marathi for to console.
\end{quote}

\subsection*{You'll want to know: \mr{हायसे} \mr{वाटणे}}
\textbf{\mr{हायसे} \mr{वाटणे}} — \emph{hāyse vāṭṇe} — \textquotedblleft{}to feel relieved\textquotedblright{}.

\begin{quote}
\mr{मुलगा} \mr{सुखरूप} \mr{घरी} \mr{आला} \mr{तेव्हा} \mr{आईला} \mr{हायसं} \mr{वाटलं}.
\end{quote}

\begin{quote}
\emph{mulgā sukhrūp gharī ālā tevhā āīlā hāysaṁ vāṭlaṁ.}
\end{quote}

\begin{quote}
Mother felt relieved when her son came home safe.
\end{quote}

Past (to him …): \mr{त्याला} \mr{हायसं} \mr{वाटलं}, \emph{tyālā hāysaṁ vāṭlaṁ}.

Dative, like \mr{वाटणे}: \mr{मला} \mr{हायसं} \mr{वाटलं}. Sentences write the neuter as \mr{हायसं}.

\begin{cousinweb}
\textbf{\mr{हायसे} \mr{वाटणे}} (\emph{hāyse vāṭṇe}) is a Marathi phrase built on \mr{हायसे}, \textquotedblleft{}eased\textquotedblright{}, whose origin is not settled, with the verb \mr{वाटणे}, \textquotedblleft{}to feel\textquotedblright{}.
\end{cousinweb}

\begin{grammarlens}[title={What you've built}]
One more word for describing what you lived and felt.
\end{grammarlens}

\subsection*{Your turn}
\begin{itemize}
\raggedright
  \item \emph{Say it:} \emph{hāyse vāṭṇe}
  \item \emph{Say it:} \emph{hāyse vāṭṇe}, once more
  \item \emph{Say it:} \emph{mulgā sukhrūp gharī ālā tevhā āīlā hāysaṁ vāṭlaṁ.} — Mother felt relieved when her son came home safe.
  \item \emph{Recall:} say the Marathi for fortitude, then the Marathi for to console, then say \emph{hāyse vāṭṇe} again
\end{itemize}

\begin{tcolorbox}[breakable,colback=teal!4,colframe=teal!35!black,title={Before you move on}]
What does \mr{हायसे} \mr{वाटणे} mean? (\textquotedblleft{}To feel relieved\textquotedblright{}.) Say it once more.
\end{tcolorbox}
\section[\texorpdfstring{\mr{गहिवरणे} (gahivarṇe)}{गहिवरणे (gahivarṇe)}]{\mr{गहिवरणे} (gahivarṇe) — to be choked up; to be overcome with emotion}

\label{lesson:MR-C458-gahivarne}



\begin{quote}
Before the new one: say the Marathi for to console, then the Marathi for to feel relieved.
\end{quote}

\subsection*{You'll want to know: \mr{गहिवरणे}}
\textbf{\mr{गहिवरणे}} — \emph{gahivarṇe} — \textquotedblleft{}to be choked up; to be overcome with emotion\textquotedblright{}.

\begin{quote}
\mr{मुलीचं} \mr{भाषण} \mr{ऐकून} \mr{वडील} \mr{गहिवरले}.
\end{quote}

\begin{quote}
\emph{mulīcaṁ bhāṣaṇ aikūn vaḍīl gahivarle.}
\end{quote}

\begin{quote}
Father was choked with emotion listening to his daughter's speech.
\end{quote}

Past (he …): \mr{तो} \mr{गहिवरला}, \emph{to gahivarlā}.

Tears near, voice thick: \mr{गहिवरून} \mr{येणे}, to be overcome.

\begin{cousinweb}
\textbf{\mr{गहिवरणे}} (\emph{gahivarṇe}) is a Marathi verb whose origin is not settled.
\end{cousinweb}

\begin{grammarlens}[title={What you've built}]
One more word for describing what you lived and felt.
\end{grammarlens}

\subsection*{Your turn}
\begin{itemize}
\raggedright
  \item \emph{Say it:} \emph{gahivarṇe}
  \item \emph{Say it:} \emph{gahivarṇe}, once more
  \item \emph{Say it:} \emph{mulīcaṁ bhāṣaṇ aikūn vaḍīl gahivarle.} — Father was choked with emotion listening to his daughter's speech.
  \item \emph{Recall:} say the Marathi for to console, then the Marathi for to feel relieved, then say \emph{gahivarṇe} again
\end{itemize}

\begin{tcolorbox}[breakable,colback=teal!4,colframe=teal!35!black,title={Before you move on}]
What does \mr{गहिवरणे} mean? (\textquotedblleft{}To be choked up; to be overcome with emotion\textquotedblright{}.) Say it once more.
\end{tcolorbox}
